\chapter{Data}\label{ch:data}

AEMO publishes the outcomes of certain variables every 5 minutes. Such quantised data is a dispatch run in its Market Management System (MMS) database. We use seven MMS tables for the year 13 September 2025 00:05 to 13 September 2026 00:00 (105,120 five-minute intervals). All are indexed by SETTLEMENTDATE, the end of the interval (AEST). They differ in data types, and in how many rows they carry per interval:
\begin{itemize}
    \item one row per region: DISPATCHPRICE (RRP, FCAS prices, the 7-day cumulative price
    and the administered-pricing flag) and DISPATCHREGIONSUM (demand, available
    generation, semi-scheduled wind and solar availability, net interchange, battery state)
    \item one row per interconnector: DISPATCHINTERCONNECTORRES (flow and import/-
    export limits)
    \item one row per binding constraint: DISPATCHCONSTRAINT, about 100 million rows over
    2,213 distinct constraints
    \item one row per locally constrained unit or blocked constraint, and one row per interval for the solver summary (CASE\_SOLUTION)
\end{itemize}

\section{Data processing}\label{sec:dataprocessing}
At the judgment of the AEMO, the dispatch engine will be overridden by manually constructed energy prices, under what is called market intervention. We only used non-interventioned data (INTERVENTION=0), because all the variables from the dispatch inform the market engine, and also because there has not been an intervention in the year interval chosen. 

Initial data filtration meant data types with no variance were discarded, as nothing can be gleaned from invariable data.  Furthermore, the constraints file reports too many constraints - 100 million rows of data simply is not feasible in a more complicated model. To resolve this over-reporting, based on the AEMO constraint FAQ \cite{aemofaq}, we dropped all \emph{non-binding} data, where the associated loss with not-adhering to the constraint was zero, keeping just those that impose rules on the optimisation of the system.  

Other features were chosen based on the descriptions given by the AEMO; With many legacy variables and heavily correlated variables, such as a duplicate of the price variable which would be unaffected by the rarely reached upper limit price, a large amount of compression was done by judgment. Because of the models chosen down the line, occasionally when a variable had a very large range in the positive direction, it was log-scaled due to the sensitivities of said model to absolute range. 